\chapter{\MakeUppercase{Technical Specification}}

\section{Requirements}
\subsection{Functional Requirements}
The system must simulate a disaster-response \ac{dtn} and execute multiple routing protocols under identical conditions. It must create heterogeneous nodes, generate messages probabilistically, detect wireless contacts, exchange messages through the selected router, expire messages after their \ac{ttl}, and compute performance metrics after each run.

The routing layer must support at least three protocols: Epidemic, Spray and Wait, and Spatio-Semantic. The proposed router must maintain encounter information, zone visit information, and role-aware forwarding decisions. It must reserve buffer space for critical messages and allow critical messages to evict non-critical messages when necessary.

\subsection{Non-Functional Requirements}
The simulator should be reproducible, simple to modify, and fast enough to run repeated experiments. The implementation should store aggregated results in \ac{csv} format so that the dashboard and thesis tables can be generated from the same evidence. The routing logic should be modular so that new protocols can be added without changing the environment.

\section{Design Constraints and Assumptions}
Several constraints shape the technical design. The project is intended for an academic thesis rather than a production deployment, so the simulator must remain understandable and inspectable. For that reason, the implementation uses explicit data structures and a centralized simulation loop instead of highly optimized event engines. This choice makes the behavior of each protocol easier to trace and compare.

The study also assumes that routing protocols should be evaluated under identical external conditions. Node population, contact range, area size, traffic generation, and \ac{ttl} must therefore remain constant across protocols within a scenario. This requirement is important because it ensures that differences in results can be attributed to routing behavior rather than to hidden changes in the environment.

Another assumption is that role-aware routing should be lightweight enough to implement in a compact simulator. The protocol does not depend on perfect map knowledge, continuous positioning infrastructure, or complex machine-learning inference. Instead, it uses information that is plausible for a field node to maintain: its own encounters, zone visits, the observed role of peers, and a small amount of transitive evidence.

\section{Feasibility Study}
The project is technically feasible because all components can be implemented using standard Python data structures and deterministic simulation steps. Contact detection uses Euclidean distance. Mobility uses a random waypoint model. Message buffers are represented as lists. Routing protocols are implemented as classes with a common exchange interface.

The simulation approach is also academically feasible because \ac{dtn} protocols are commonly evaluated through simulation due to the cost and variability of field deployments. Tools such as the ONE simulator demonstrate the established practice of simulation-based \ac{dtn} protocol evaluation \citep{keranen2009one}. This project uses a custom simulator so that node roles, critical-message logic, and protocol internals can be controlled directly.

\section{System Specification}
The major simulation parameters are shown in Table \ref{tab:sim_parameters}. These parameters intentionally create a stressed environment: a large area, short contact range, finite buffer, high critical-message proportion, and multiple traffic levels.

\begin{table}[h!]
	\centering
	\caption{Simulation parameters}
	\renewcommand{\arraystretch}{1.3}
	\resizebox{\textwidth}{!}{
	\begin{tabular}{|l|l|l|}
		\hline
		\textbf{Parameter} & \textbf{Value} & \textbf{Purpose} \\
		\hline
		Simulation area & 2200 by 2200 units & Creates sparse contacts and movement diversity \\
		\hline
		Transmission range & 90 units & Limits contact opportunities \\
		\hline
		Buffer size & 30 messages per node & Forces meaningful buffer decisions \\
		\hline
		Simulation duration & 8000 time steps & Allows encounter and zone history to develop \\
		\hline
		Message \ac{ttl} & 2500 time steps & Penalizes delayed or wasted forwarding \\
		\hline
		Critical-message probability & 0.45 & Stresses priority handling \\
		\hline
		Critical copies & 16 initial copies & Allows urgent messages to spread when useful \\
		\hline
		Non-critical copies & 8 initial copies & Provides controlled replication for ordinary messages \\
		\hline
		Traffic levels & Low, Medium, High & Tests behavior from light load to congestion \\
		\hline
	\end{tabular}}
	\label{tab:sim_parameters}
\end{table}

\section{Parameter Selection Rationale}
The selected parameters are intended to create a meaningful routing challenge instead of an artificially easy environment. A large map and a relatively short radio range reduce the probability of constant connectivity. This encourages store-carry-forward behavior and makes carrier choice important. If the area were smaller or the range much larger, most protocols would appear similarly strong because frequent contact would mask their differences.

Finite buffers and a moderate \ac{ttl} are equally important. If buffers were effectively unlimited, flooding would face little penalty and the value of selective forwarding would be harder to demonstrate. If the \ac{ttl} were too long, poor routing choices could still succeed eventually, reducing the practical significance of delay. The chosen values therefore create a balance where both delivery and timeliness matter.

The critical-message probability and copy allocation are also deliberate. A high critical-message proportion forces the router to make priority decisions often enough for the effect to be visible. The larger initial copy count for critical messages reflects the thesis assumption that urgent communication deserves additional forwarding opportunity, provided that replication remains controlled by utility-guided logic.

\section{Node and Traffic Model}
The network contains 35 civilians, 8 responders, 3 shelters, and 5 drones. Civilians move slowly, responders move at medium speed, shelters remain static, and drones move quickly. The shelter nodes are placed in a clustered region to represent emergency coordination points. Drones are initialized as fast relays and are useful for carrying messages across partitions.

Messages are generated at each time step according to the selected traffic probability. A generated message chooses a random destination other than the source. Each message has a creation time, unique identifier, hop count, critical flag, copy count, and destination. Delivered messages are counted once and then removed from all buffers.

\section{Requirement-to-Metric Traceability}
The technical requirements are connected directly to the evaluation metrics. The requirement to preserve overall communication success is measured through delivery ratio. The requirement to protect urgent traffic is measured through critical delivery ratio and average critical delay. The requirement to avoid uncontrolled replication is reflected in overhead ratio. The requirement to manage limited storage is reflected in buffer drops.

This traceability is useful because it links the technical specification to the final discussion. A protocol cannot be described as successful merely because one metric improves. For example, higher delivery achieved at the cost of severe congestion would not satisfy the non-functional goal of resource efficiency. Similarly, low overhead alone would not satisfy the functional goal of timely emergency delivery. The metric set therefore acts as a structured interpretation framework for the results chapter.

\section{Scalability and Extension Considerations}
Although the simulator is compact, the technical specification is written so that the system can be extended. The modular routing interface means that additional \ac{dtn} protocols can be inserted into the same environment with limited code change. Likewise, the node-role model can be expanded to include more specialized actors without rewriting the full simulation loop. This extensibility is part of the project quality requirement because a thesis implementation should support further experimentation rather than end with a single fixed comparison.

Scalability in this thesis does not primarily refer to very large distributed deployment. Instead, it refers to whether the simulation framework and routing logic remain manageable as the study grows in scope. By keeping the maintained state compact and by storing results in simple \ac{csv} form, the project supports repeated runs, added metrics, and additional scenarios without requiring a major redesign of the codebase.

\section{Performance Metrics}
The following metrics are used:

\begin{itemize}
	\item Delivery ratio: delivered messages divided by generated messages.
	\item Critical delivery ratio: delivered critical messages divided by generated critical messages.
	\item Average delay: average time between message creation and delivery.
	\item Average critical delay: average delay for critical messages only.
	\item Overhead ratio: transmissions divided by delivered messages.
	\item Buffer drops: number of insertion failures or forced drops caused by buffer pressure.
\end{itemize}

These metrics reveal different protocol properties. Delivery ratio measures reachability, critical delivery ratio measures emergency effectiveness, delay measures timeliness, overhead ratio measures resource cost, and buffer drops measure congestion pressure.
